\documentclass[11pt]{article}
\usepackage{../shared/luxcover}
\usepackage{amsmath,amssymb}
\usepackage{hyperref}
\usepackage{booktabs}
\usepackage{enumitem}
\usepackage{xcolor}

\title{Lux Exchange}
\author{Zach Kelling\\\textit{Lux Industries, Inc.}\\\texttt{research@lux.network}}
\date{February 2023}

\begin{document}

\luxcoverpage

\begin{abstract}
``I believe that banking institutions are more dangerous to our liberties than standing armies. If the American people ever allow private banks to control the issue of their currency, first by inflation, then by deflation, the banks and corporations that will grow up around will deprive the people of all property until their children wake-up homeless on the continent their fathers conquered. The issuing power should be taken from the banks and restored to the people, to whom it properly belongs.''

-- Thomas Jefferson
\end{abstract}

\tableofcontents
\newpage


-- Thomas Jefferson


``One of the greatest gifts to humankind is the ability to create technology. We have been granted the tools to improve every facet of our collective existence. As a civilization united under one purpose: the betterment of humankind through the enrichment of the human life cycles of trade. We, The Lux community, dedicate the application of this technology to provide comfort and happiness with a sustainable and symbiotic relationship to the planet earth.''


-- The Lux Community


\section{Executive Summary}

This White Paper describes the core concepts behind Lux Exchange (LX), a novel solution to solving the front running problem of other decentralized exchange platforms. The Lux Exchange is a high performance and fully  Interoperable Social Trading DEX (Decentralized Exchange) built on the Lux Network and secured with a zero knowledge bridge, Z-Chain.


LX main features are:


Social Features: Trading digital assets on LX is not only intuitive but also easily allows users to copy trade, mirror trade, and build people based portfolios. Users can create profiles, post media content, and interact with other users on the exchange platform.


Reliability: LX operates on the Lux Network. As shown in Figure 1, at any given time a system node is designated as Leader to generate a Proof of History sequence, providing the network global read consistency and a verifiable passage of time. The Leader sequences user messages and orders them such that they can be efficiently processed by other nodes in the system, maximizing throughput. It executes the transactions on the current state that is stored in RAM and publishes the transactions and a signature of the final state to the replication nodes called Verifiers. Verifiers execute the same transactions on their copies of the state, and publish their computed signatures of the state as confirmations. The published confirmations serve as votes for the consensus algorithm. Figure 1: Transaction flow throughout the network. In a non-partitioned state, at any given time, there is one Leader in the network. Each Verifier node has the same hardware capabilities as a Leader and can be elected as a Leader, this is done through PoS based elections. In terms of CAP theorem, Consistency is almost always picked over Availability in an event of a Partition.


High Liquidity: Liquidity Aggregator Smart contracts forked from Uniswap V3 allow Platform traders the ability to provide marginal liquidity mining for an increased yield directly proportional to the limits assessed in the yield farming option.  This increases the potential gains for the liquidity provider and also creates more incentive in making LX highly liquid.


High Performance: The core platform is built with industry standard general purpose programming language (C++) and a binary instruction format for stack-based virtual machines (solidity/hardhat scaffold infrastructure compiled into Solidity and executed on the Lux Network). These modern technologies enable enterprise level transaction speeds.


Simple Trading: a simple portfolio management interface is provided for traders.  Users are aligned with their target market. Users don't have to manage multiple trading accounts with time sensitive operations. Firesale Feature allows Traders to quickly exit all of their positions quickly with one click of the Firesale button.


Diversity: diversity is essential for lowering risk when trading stocks, crypto-currencies, commodities, and digital assets. Users can spread out their invested capital across various markets to diversify their portfolios.


Interoperability/Bridge: A blockchain bridge is a link that allows the transfer of tokens and/or arbitrary data from one blockchain to another. Both blockchains can have different protocols, rules and governance models, with the bridge it provides a compatible way to interact securely on both sides. Through advanced multi token pairs and wrapped token chain transfers, LX provides full interoperability amongst multiple blockchains and allows many different blockchain ecosystems to interact on the LX order books through oracles* for decentralized trading.


User Support: LX user support provides a live chat experience. The Lux Curated Registry Customer service model provides support at all levels of the trading experience. This helps advanced traders who need help writing code for trading bots, and beginners who need help navigating the site or putting together their first portfolio.


Transparency: User profiles contain a transparent outline of their portfolio holdings, pre-calculated risk assessment rating, and statistical record.


Security: Falling victim to a brute force attack is always a possibility as the processing power of commodity computing systems continues to improve. LX uses state-of-the-art commercially available cryptographic and hashing methods to achieve the smallest possible brute force attack surface. That includes the use of proven key management techniques as well as homomorphic encrypted zero knowledge file systems (Lux Zk Bridge).


*Oracles are computer programs that connect data from the outside world (off-chain) with the blockchain world (on-chain). Decentralized oracles achieve trustlessness and deterministic results that rely on cause- and-effect rather than on individual relationships. By leveraging many different data sources, and implementing an oracle system that isn't controlled by a single entity, decentralized oracle networks provide an increased level of security and fairness to smart contract order books. Lux Exchange Implements Chain.Link for all of the offchain data that it is interacting with.


\subsection{Negative Consequences of Centralization}

Decentralization and distribution of an immutable ledger are the key concepts that make Blockchain technology so incredibly valuable. Yet, ironically, most cryptocurrency exchange platforms are centralized. The security of blockchain technology doesn't exist in the architecture of centralized exchange platforms. Centralization introduces the Single Point of Failure (SPoF) problem into the system. More specifically, centralized databases and ledgers within the architecture of centralized platforms make the whole system vulnerable to both digital and physical attacks.


The meaningful contribution of a centralized exchange system is that they possess the power to control the outcome of a transaction in the advent of a financial  mistake or error within the system's internal programming. This leads users to feel as though that central body would be held accountable for the loss of one's digital assets and be able to correct that mistake. However, time and time again, centralized exchange platforms fail to provide clear accountability for assets stolen from their users. Answers to questions regarding security protocols, safety of digital assets and customer service are usually obscure, if not nonexistent. Centralized exchange platforms hold user keys and control the leverage and movement of their funds.


Centralized exchange platforms often store funds in online hot wallets and many have failed to protect the private keys of user wallets, which is why they have fallen victim to many hackers. More than 980,000 BTC have been stolen from these centralized exchanges (nearly \$40 billion USD at current rates). Few of these exchanges were insured or regulated. Centralization in any system introduces security vulnerabilities and on centralized exchange platforms can accommodate fraudulent behavior. The following list provides information about some of these events:


Silk Road


When: October 2013


Amount stolen: BTC 171,955 (USD \textasciitilde{}270,000,000)


Comments: Silk Road was a marketplace that accepted cryptocurrencies, not an exchange. Still, users had BTC linked to their accounts, and keys were stored in a centralized manner. After the FBI took Silk Road down, all BTC was confiscated and users lost their assets.


MtGox


When: March 2014


Amount stolen: BTC 850,000 (USD \textasciitilde{}700,000,000)


Comments: The biggest hack in the history of cryptocurrencies, with one of the most popular exchanges at the time. MtGox failed to protect the private keys of user wallets, and hackers drained funds away. The amount of stolen BTC is currently worth more than 6 billion USD.


Cryptsy


When: January 2016


Amount stolen: BTC 13,000 + LTC 300,00 (USD \textasciitilde{}9,500,000)


Comments: A Trojan malware was inserted into Cryptsy's code, and the attacker was able to transfer BTC and LTC out of the exchange's wallets.


Mintpal


When: December 2014


Amount stolen: BTC 3,894 (USD \textasciitilde{}3,200,000)


Comments: Mintpal was one of the most popular trading platforms. Customers were told the company was going to have new ownership. The new owner failed to detect and patch the vulnerabilities that led to the platform being hacked, and many believe the acquaintance and the hack were nothing but an inside job.


Bitstamp


When: January 2015


Amount stolen: BTC 19,000 (USD \textasciitilde{}5,100,000)


Comments: Bitstamp employees received a malicious file that infected the company's infrastructure.


Bter


When: February 2015


Amount stolen: BTC 7,000 (USD \textasciitilde{}1,750,000)


Comments: Bter had been previously hacked, so it was not the first incident in this centralized exchange.


Bitfinex


When: August 2016


Amount stolen: BTC 120,000 (USD \textasciitilde{}72,000,000)


Comments: Bitfinex is famous for the creation of Tether and sharing executives. They advertised as having multisig wallets for each customer, which still did not prevent them from losing users BTC.


NiceHash


When: December 2017


Amount stolen: BTC 4,000 (USD \textasciitilde{}60,000,000)


Comments: NiceHash is not an exchange but a cloud mining service. In December 2017 their servers were hacked and miner wallets were emptied.


Coincheck


When: January 2018


Amount stolen: NEM 523,000,000 (USD \textasciitilde{}534,800,000)


Comments: Coincheck used cold wallets for its BTC trading operations, but neglected security measures on Asian currency NEM. All NEM deposits were stored in a single wallet that was emptied during the hack. This could have been avoided by securing hot wallets with greater levels of encryption standards.


BitGrail


When: February 2018


Amount stolen: NANO 17,000,000 (USD \textasciitilde{}195,000,000)


Comments: BitGrail failed to secure it's new 0-fee cryptocurrency NANO storage. NANO used the recently introduced block lattice, and the fact that a centralized exchange was adopting such infant and unvalidated technology might have influenced this episode.


CoinSecure


When: April 2018


Amount stolen: BTC 438 (USD \textasciitilde{}3,300,000)


Comments: After the hack, CoinSecure's owners filed a lawsuit against one of the exchange's employees claiming it was an inside job.


\subsection{1 Security Solutions as a Decentralized Exchange}

Operating as a Decentralized Exchange (DEX) platform, LX addresses the problems of security vulnerabilities in centralized exchange models with private key storage and performance issues in current DEX platforms. This is achieved by securing the database within the Lux Network Infrastructure. The following strategies are implemented for maximum platform security and performance:


Critical Core functionality (e.g.: smart contracts, matching engine) is compiled into binary executables increasing computational speeds. Scripting languages (e.g.: JavaScript, Python) are restricted to non critical functionality (e.g.: UX, Social Features).


User Interface is enforced by mandatory multi factor (MFA) authentication.


SHA3, a quantum resistant hashing algorithm is used to secure and encrypt authenticated user data.


Industry proven Symmetric Cryptographic Algorithms protect network traffic.https://eprint.iacr.org/2018/379.pdf


Industry proven Public Key Infrastructure (PKI) algorithms authenticate users and blockchain nodes.


Assets are stored in Time-Delayed, Multisignature controlled accounts.


Benign Fill Techniques ensure key material is never stored or transmitted in an unprotected state and is never in the hands of humans in cleartext form.


Multisignature is used on hot wallets with several independent processes/servers double-checking all withdrawals.


Multisignature is required to whitelist accounts.


Hardware Wallets are required for all signing, even automated withdrawals.


\section{LX: A Decentralized Social Trading Platform}

\subsection{1. Overview}

LX is a high performance DAO (Decentralized Autonomous Organization) operating as an open source social trading platform. It is designed for peer-to-peer exchange of cryptocurrencies, digital assets, market pegged assets, stocks, and commodities. Users can build and manage portfolios using copy trading, mirror trading, and people based portfolios, while interacting with other users on the platform.


Smart contracts handle all the financial transactions using actions, handlers, context free actions and multi-index containers on the Lux Network. LX's smart contracts are implemented as Free/Libre Open Source Software for easy audit of the source code. The code for the smart contract is made available under a public repository.


\subsection{2. Decentralized Application}

The core of the LX Decentralized Application (DApp) is built with the industry standard general purpose programming language (C++). It is compiled into Solidity binary executables, and operates on the Lux  Blockchain.


Amongst the LX DApp security features:


Signed Transactions: Lux Wallet enables users to sign transactions with their private keys securely from their phone or browser without ever exposing them. LX uses Lux Wallet to provide a secure channel of communication between the browser and Lux Network.


Anonymity: the user of the average exchange must trust that their IP information is not being recorded. They also must trust that the data presented to them is correct and matches with what is on the blockchain. LX Dapp uses smart contracts with permission mapping to fetch data directly from the blockchain and store it locally at browser runtime, eliminating the need for a trusted third party. This decentralized approach results in improved anonymity and trust.


Hardware Wallet support: the LX DApp enables support for cold storage hardware wallets such as Trezor and Ledger and is not compatible with wallets that do not meet the security requirements set as digital asset safety standards.


Reduced Attack Surface: nobody is immune to large-scale internet attacks. Decentralized Applications and Blockchain technology directly mitigate many attack vectors by design.


Quality Assurance: all code for the LX platform complies to IEC 61508 and CERT Coding Standards.


The LX DApp User Interface (UI) is written in TypeScript. The UI code makes explicit calls to the Solidity binaries for execution of smart contracts and core functionality inside the browser runtime engine. Solidity is the state-of-art web technology promoted by the World Wide Web Consortium (W3C) W3C technology, which means it is supported by all major modern browsers in the market.


User experience is a key success factor for any exchange. The LX DApp provides cross platform mobile and web interfaces to make sure users are able to trade, issue new user-defined assets, receive market information, as well as potentially making payments with some of its price-stable assets.


All code for the LX platform is available on public repositories for public auditability and community support.


\subsection{3. User Experience}

For a novice trader, cryptocurrency trading platforms tend to offer a sterile and somewhat complicated User Experience (UX). The LX Social Exchange Platform User Interface is designed solely with the first-time trader in mind. Back in 2006, Jakob Nielsen published an article about the "F-Shaped Pattern For Reading Web Content". Nielsen's team found that users scanned the web content in the so-called F-shaped pattern. In an eye tracking study 232 users were recorded navigating thousands of web pages, all showing the same patterns. According to the article: "Users first read in a horizontal movement, then move down the page a bit, then read across and finally scan the content's left side in a vertical movement. This creates an F - shaped reading pattern (which can sometimes look like an E). Social media websites have utilized such patterns and now we bring this facebook-like UX to trading.


\subsubsection{Accounts, Wallets, and Keys}

Each user is uniquely identified by their account. Wallets are used to store assets. Each wallet is uniquely tied to a key pair, where the public key represents the wallet's address, and the private key functions as proof of wallet ownership.


\subsubsection{Authentication}

Users gain access to their accounts through authentication. Each wallet is locked and unlocked with a unique high entropy password that users keep as a secret. With ecosystem partner LUX DATA-SECURITY , Biometrics authentication and 4FA will ensure all keys are secured to the highest available standards.


\subsubsection{Hardware Wallets}

Hardware wallets are portable devices that safely store private keys. Tamper proof chips erase internal memory if microscopic seals are broken, which keeps keys from ever being seen by human eyes. Communication with the DApp happens via the Chip Card Interface Device (CCID) USB protocol.


LX allows users to choose between storing their private keys in their own hardware or inside the Yubikey NEO, a hardware authentication device manufactured by Yubico that supports one-time passwords, public-key encryption and authentication, and the FIDO Alliance's Universal 2nd Factor (U2F) protocol.


Yubico's GitHub repository provides low level C libraries for hardware integration, a U2F Server C Library, amongst other resources for the LX DApp development.


\subsubsection{Portfolios}

A portfolio is a collection of wallets tied to a user's account. Each wallet represents proof of ownership of a different kind of asset. Assets might be stocks, crypto-currencies, commodities, or digital assets.


The Portfolio Management Tool allows users to organize digital assets. Corporate actions like dividends and stock splits are automatically incorporated. Users can analyze their past actions and learn from past successes and errors.


\subsubsection{Social Trading}

Social trading is a service that allows investors to replicate the operational movement of expert traders. Social trading requires little to no knowledge about financial markets and users only need to make a correct selection of traders to follow and learn from. This assessment is available through transparent financial statistics that provide insight about the user's successful trades and their diversification of financial instruments.


Social trading provides a model to compare and copy trades, techniques and strategies. Prior to the advent of social trading, investors and traders relied on fundamental technical analysis to form their investment decisions. With social trading, investors and traders can integrate social indicators from other traders into their investment decision-process.


Social trading has been described as a way to shorten the learning curve from novice to experienced Forex trader [?]. Traders can interact, watch others take trades, then duplicate their trades and learn what prompted the top performer to take a trade in the first place. By copying trades, traders can learn which strategies work and which do not work [?].


LX provides the following social trading features for its users:


A social profile with username, avatar, approximate regional location. This profile represents the user's account.


Trade history data: user's success rate, user success ranking, average daily profit/loss, annual gains percentage.


User friends and friend requests.


Trader risk rating, defined by system analysis.


Community position rating.


Successful investment ranking.


\subsubsection{People Based Portfolios}

People Based Portfolios (PBP) are organized sets of transaction patterns. Traders can choose to create a PBP and share their trading activity with the community. Novice users can analyze successful PBPs and learn from successful trading patterns. Successful traders are rewarded LUX proportionally to how many users are copying their trades.


\subsubsection{Copy Swaps}

Copyswap protocol is a decentralized asset management protocol that allows expert traders the ability to create permission mapped liquidity aggregator smart contracts. These novel and unique smart contracts serve to facilitate the decentralized and trustless management of digital assets.  When an investor finds a particular portfolio that is performing in a way they deem desirable for investing into, they can assign a percentage  of their digital assets to go under management while remaining full custodial of their keys. All assets remain on chain while assigning permission mapping of asset management to allow expert traders to move investor money in a way that is trustless and decentralized, mitigating against potential risk of fraud like ponzi schemes like in the case of Enron or Bernie Madoff. Users can invest their digital assets into the protocol with zero fees until there are gains realized while those digital assets that are entered into the Swap liquidity contract.


\subsubsection{Trading Charts}

Trading charts continuously inform users about the decentralized liquidity pool orderbook. The charts plot technical indicators such as closing price, trading volume and moving averages.  This helps identify market trends and guide trading strategies.


\subsubsection{Indicator Alarm Manager}

The Indicator Alarm Manager (IAM) notifies traders about important events. Users can receive alerts about interaction between moving averages across time, critical economic data updates, as well as and when specific values are breached in indicators (e.g.: RSI14).


\subsubsection{Smart Search}

LX provides a discovery tool that fully indexes content from keywords, traders, news and databases in a single search box.  It also allows users to search order books, pairs, and collections as well as community members and @token brands to identify community engagement socially as well as financially with those digital assets.


\subsubsection{Watchlist}

LX provides a dashboard that shows trading data from all LX markets. Traders can explore and analyze market data easily.


\subsection{4. Community Support}

The LX social exchange platform incentivizes community activity in two major ways: Rewarded Content Production and Token Curated Customer Service.


\subsubsection{Rewarded Content Production/Trading Bots}

Developers and traders are incentivized to produce and curate content inside the LX platform in a similar fashion to Steemit.com. Users can interact in the LX network to discuss crypto and stock markets and share trading tips on commodities. Content quality is proportionally rewarded by community up and down voting.


Trading Bots: An anonymous 0-fee exchange is an ideal environment for trading bots. LX fundamentally believes that in an unregulated market that works 24/7, the best path towards price stability is market-pegged assets and widespread use of trading algorithms that exhibit market inefficiencies. LX allows implementation of trading bots by developers who have the knowledge of how markets operate, as well as for traders who know how to code and want to get into programmatic cryptocurrency trading.


Examples of trading bots include, but are not limited to:


Market making bot that provides liquidity for any token-token trading pair.


Arbitrage with other decentralized exchanges to make sure LX users have the best price.


Media marker indicators.


Automated Fibonacci retracement.


Amongst the tools provided for bot writers is strategy backtesting and download of historical order book and trade data. Developers are able to use the LUX API to write DApps, trading bots, or anything related to LX and LUX. This programmatic content can be shared in a similar fashion to Utopian.io.  In order to incentivize quality assurance of source code, developers are also rewarded for auditing other developer's code.


\subsubsection{Token Curated Customer Service}

Customer support and user education are crucial in maintaining a trusted and efficient user experience. The primary support system of the LX exchange is a live chat interface, connecting the user with a support team member 24/7. Users are incentivized to become support providers by receiving LUX Tokens for their services. They are required to stake tokens and pass service representative training courses to earn a customer rated token income for quality performance.


The ranking system is automated and maintained by an immutable token curated registry of multiple levels. As a customer service representative earns higher ranking, they earn more tokens for their time during phone calls and chats. The incentive to provide quality support is the staked tokens each representative is required to put as collateral at each rank upgrade. As the number of customer interactions increases, the customer service representative is incentivized with an opportunity to apply for the next ranking test.


\subsection{5. LX Architecture Comparison}

This section demonstrates how LX's architecture stacks up against other centralized and decentralized exchanges.


\subsubsection{eToro}

eToro is the multi asset brokerage company that was the first to introduce social trading on a mass scale. The eToro OpenBook enables investors to view, follow and copy the network's top traders automatically.


The main problem with eToro is centralization. eToro's system is not asynchronous BFT. If bad actors are able to control eToro's servers and order books, the whole system is put at risk.


\subsubsection{EtherDelta}

EtherDelta is an exchange with hybrid centralized and decentralized features. Most operations are handled directly on the blockchain. A trader never needs to let go of the the ownership of their tokens, enabling the system to deal with transfers and trades in a trustless manner.


The weak point of EtherDelta's architecture is the centralized order book server. In case a bad actor took control of the orders, they would be able to trade against those orders or delete them. Deletion of orders would remove all liquidity from the exchange and render it unusable.


\subsubsection{0x Protocol}

0x is an open, permissionless protocol allowing for ERC20 tokens to be traded on the Ethereum blockchain. 0x protocol is a pluggable building block for DApps that require exchange functionality on the Ethereum blockchain. The 0x protocol mitigates the security issues faced by centralized exchanges by enabling fast decentralized transactions.


0x uses state channels to move transactions off the main blockchain by making participants send signed messages between each other. After the microtransaction is complete, the result is published to the main blockchain and gas fees are applied.


The main drawback of 0x protocol's state channels is the need for participants to be online during state channel transaction, security deposits, and the many on-chain transactions required to set up the state channel. This makes it inefficient for one-time transactions.


Another issue with 0x is the fact that the Ethereum network transaction speeds make it not ideal for decentralized exchanges. The high latencies often interfere with price fluctuations, forcing traders to execute orders that they would have opted out of were latencies lower.


\subsubsection{LX}

LX combines the security of zero knowledge homomorphic encryption with distributed ledger technology to provide a novel and secure asset exchange protocol that is trustless, decentralized, and private as it executes shielded transactions to automate its Defi smart contracts.


The use of Lux permission mapping and liquidity aggregators provide an additional unique and novel feature called copy trading. Copy trading is provided through a DAO operated liquidity pool that allows users to combine both fungible and non fungible assets into a STF (single trade fund).


Concentrated liquidity, giving individual LPs granular control over what price ranges their capital is allocated to. Individual positions are aggregated together into a single pool, forming one combined curve for users to trade against.


\section{Lux Protocol}

Lux protocol combines all the benefits of Uniswap V3 with new multi-asset LP improvements. With Lux Protocol, LPs can stake both fungible and non fungible digital assets. This new feature provides automated market making of NFT's that allows them to be  swappable in exchange for other NFT's, approximating a fee-earning limit order that executes along a smooth curve.


LPs can provide liquidity with up to 4000x capital efficiency relative to Uniswap v2, earning higher returns on their capital.


Capital efficiency paves the way for low-slippage trade execution that can surpass both centralized exchanges and stablecoin-focused AMMs


LPs can significantly increase their exposure to preferred assets and reduce their downside risk


LPs can sell one asset for another by adding liquidity to a price range entirely above or below the market price, approximating a fee-earning limit order that executes along a smooth curve


\subsection{1. Distributed Autonomous Organization (DAO)}

The LX platform operates as a Distributed Autonomous Organization (DAO). The concept of a DAO arises as an autonomous organization with a set of self executable rules implemented as smart contracts whose guidelines/rules are set by stakeholder voting.


These strict rules are implemented as publicly auditable Free/Libre Open Source Software distributed across stakeholder computers. By owning the LUX Token, stakeholders are entitled to vote on protocol upgrades.


\subsection{2. Governance}

As Bitcoin and Ethereum have already demonstrated, even in the best decentralized projects, developers hold more power over the project than any other stakeholders. In case of Bitcoin, the developers of the blockchain were the main reason for the Bitcoin Cash and SegWit2x hard forks. In the case of upgrades so, the developers chose to create the Ethereum Classic fork after the DAO hack. When the power of a decentralized system is in the hands of a small group of people, the trust in the protocol is undermined and people are less likely to embrace it.


The Lux Protocol introduces voter permissions, which makes the decision making a more democratic process. The voting is tamper-proof and conducted on the Lux Network. Every LUX holder can vote for changes to the Lux Protocol and participate in its development. The price of LUX on the market act as a measure of the protocol's importance, as only LUX holders are allowed to apply for a proof of humanity token to participate in voting. This ensures that LUX holders are incentivized to keep the fees as low as possible in order to maintain market liquidity and trade volume. The more trades that happen through the Lux Protocol, the more power LUX holders gain, and the greater the market value of the LUX gets.


The Lux Protocol expands the concept of shared infrastructure ownership. Power imbalances between stakeholders are mitigated by the introduction of democratic governance mechanisms through smart contracts. Smart contracts are publicly auditable set of predefined rules with power over entities and funds. LUX Token holders can vote on almost every aspect of the Lux Protocol. All network parameters, from fee schedules to block intervals and transaction sizes, can be tuned via community voting. Users can even vote for others to vote in their stead. Proxy voting ensures that everyone in the ecosystem is properly represented, even if they do not have the time or the inclination to weigh in on every single issue.


LX's DApp allows LUX holders to navigate votes, discuss and coordinate proposals, and vote for changes in the LUX. The DApp design prioritizes ease of use and accessibility, making governance of the exchange platform fair and democratic.


\subsection{3 LX Tokenomics}

In this section Lux exchange will outline the mathematical equations chosen to calculate the staking and rewards models as well as the algorithms that are used to assess the value of the governance token and liquidity pool tokens.


\subsubsection{3.1 Equation of exchange}

The equation of exchange is defined as: MV=PT


Where: M= money supply, V= velocity of money, P= average price level of goods, T= index of expenditures (such as the total number of economic transactions)


In token economies, this has been adopted by two prominent people --- Chris Burniske and Vitalik Buterin.


Burniske definition: MV=PQ


Where: M= size of the asset base, V= velocity of the asset (the number of times that an average coin changes hands every day), P= price of the digital resource being provisioned, Q= quantity of the digital resource being provisioned


Using the Burniske definition, valuations typically solve for M by rearranging the equation: M=PQ/V


In order to solve for token price, one must calculate M, by working out the size of the market in dollars (PQ), divide it by the velocity (V) and then divide M by the number of coins in supply.


Buterin definition: MC=TH


Where: M= total money supply (or total number of coins), C= price of the currency (or 1/P, with P being price level), T= transaction volume (the economic value of transactions per time), H= 1/V (the time that a user holds a coin before using it to make a transaction)


Using the Buterin definition, to solve for the token price, one must solve for C:


C=TH/M


In either definition, one can see that the velocity of the coin is inversely proportional to the value of the token i.e the longer people hold the token for, the higher the price of each token. This is intuitive, because if the transactional activity of an economy is \$100 billion (for the year) and coins circulate 10 times each over the course of the year, then the collective value of the coins is \$10 billion. If they circulate 100 times, then the collective coins are worth \$1 billion. Thus, understanding and calculating the velocity in any token economy is extremely important.


\subsection{LUX Exchange Token Staking and Rewards}

How to Stake LX?


Staking your LX tokens on Lux Exchange is the best way you can help secure the world's highest-performing blockchain network, and earn rewards for doing so!


In order to stake tokens on Lux Exchange, you first will need to transfer some LX into a wallet that supports staking, then follow the steps or instructions provided by the wallet to create a stake account and delegate your stake.


The whole process might take you not longer than 20minutes, ànd you will keep full control of your funds.


There is no required minimum amount to stake and the rewards are automatically reinvested.


For more information on Staking LX, please read our Step-by-Step Guide on How to Stake Lux Exchange.


How much can I make staking LX?


Lux Exchange has a proposed 8\% inflation rate which is gradually decreasing by 15\% per year, until reaching a floor of 1.5\%.


The proposed inflation is considering a 400ms slot time. However the current average slot time could last longer, and thus the inflation is considerably lower.


The inflation is distributed together with transaction fees between all validators proportionally to their staking balance.


Validators share the rewards with their delegators after deducting their commission.


To estimate your Staking Rewards check out the Lux Exchange Staking Rewards Calculator here.


How to choose the right Lux Exchange Validator?


When choosing a validator to delegate your LX tokens, make sure to look at the commission rate which has a direct impact on your rewards being paid out. Low Commission = High Rewards.


However please note that commission rates are the bread and butter for many validators. They rely on it to maintain reliable operations.


Furthermore, it is important to make sure that the selected validator is able to maintain a solid close to 100\% uptime for their services.


You may want to consider delegating to smaller validators in order to further decentralize the network. This does not only support the network resilience, but also the value of your LX investment long-term.


Also consider validators that are long-term committed to providing value to Lux Exchange by supporting the platform's app development, tooling, and educational materials.


At Staking Rewards we have pre-vetted a bunch of validators in Lux Exchange. Browse through the list above and search for the verified badge. These validators are registered with Staking Rewards and are considered reliable. Please note that any validator is free to get in touch with us to get verified. We are not affiliated with any of them and thrive to provide an independent ranking.


Is there any risk to staking LX?


There is no significant risk when delegating LX.


Please consider that there is a risk of slashing up to 100\% of your funds, in case that the validator you delegated to, signs illegal transactions or votes for illegal forks. However if the validator and your funds get slashed would be decided case by case based on an on-chain governance vote.


If you choose a reliable validator, the risk of slashing should be close to 0.


Additionally, please consider that there is a 2 days un-bonding period before you can access your delegated funds and transfer them.


\subsection{4. Decentralized Liquidity Pool (DLP)}

The Decentralized Liquidity Pool (DLP) is a set of liquidity aggregator smart contracts that aggregates trades across a multitude of exchange platforms and is, essentially, a decentralized and distributed order book. Creation of orders and trades against existing orders occur inside the DLP. Order actions assess multiple market asset prices to provide the quickest and most efficient way of matching trades. This allows DApps and other exchange platforms to utilize this liquidity pool by staking pairs of tokens for an incentive without having to ask permission. Ecosystem benefits include that Lux.network dApps get greater liquidity, and the Lux Protocol benefits from facilitating a larger trade volume.


The DLP order book is stored directly on the LX subnetwork. In the advent that by some catastrophe, the LX block producers were to become compromised by bad actors, the DLP order book could be recreated from scratch. This is done by reading all events, starting from the block when LX was deployed on the blockchain and recreating the order book based on this information. This allows independent exchange implementations to be in sync among each other without having to trust a centralized API maintained by an exchange. In other words, the Lux Protocol treats the blockchain as an append-only log. LUX builds on the architecture that is decades old in database research.


\subsubsection{Market Maker Fees}

Being adaptable is the key to success for a decentralized exchange protocol. Instead of having fixed fees, it is beneficial to retain certain flexibility in the following scenarios:


If the protocol is getting abused by wash traders on some particular token, LX Token holders may vote to raise the fees in order to discourage such behavior.


When the protocol is fighting for user preference with other decentralized exchanges, LX Token holders may vote for lowering the fees.


When the Lux Protocol needs to be upgraded, LX Token holders may vote to take the designated portion of market maker rebates and pay a team of developers to improve the protocol.


When the protocol has firmly established its place in the token ecosystem, LX Token holders may vote to remove the fees altogether.


In the future it might make sense to add new liquidity pools for other token standards, such as non-fungible tokens. In this case, a new DAO smart contract would be created, and the DAO could vote to transfer the Lux Protocol ownership to the new smart contract. This would allow the DAO to adopt new features of the blockchain while keeping the same LX Token, contributing to the token's long-term value.


\subsection{DAO Proposals for Protocol Upgrades}

Each proposal would require tokens to be staked in order to prevent a spam attack on the DAO, and there would be a minimal amount of LUX required to present a proposal to be voted on. Each proposal would be active for a maximum of 30 days. Proposals would require staking user funds. As soon as a proposal passes the success threshold, its desired effect happens immediately via the power of smart contracts via Lux  blockchain protocol execution of context free actions. The staked tokens from this vote would be available for withdrawal, and the changes to Lux Protocol would be performed automatically as part of the last voting transaction. Each token can only be used once per vote and is available for withdrawal after the vote is recorded onto the Lux Network.  If the vote failed to pass the success threshold for 30 days it would be considered to be a failed vote. No changes to Lux Protocol would be made, and those who had staked their tokens in this vote would not be eligible to make a percentage gain of the successful proposal outcome model. The successful proposal model is a reflection of the positive outcome for the use case described in the ricardian contract.


\subsection{4. LX C++ Application Programming Interface (API)}

The LX C++ Application Programming Interface (API) is the main way DApps can interact with the Lux Protocol. Amongst the main functions provided by the LUX API are:


Deposit


Check Balance


Withdraw


Create Order


Cancel Order


Execute Order


Get Order Book


\subsection{5. Permission Mapping}

The Lux  platform allows each account to define a mapping between an action or contract of any other account and their own named permission level. For example, an account holder could map the account holder's social media application to the account holder's ``friends'' permission group. With this mapping, any friend could post as the account holder on the account holder's social media. Even though they would post as the account holder, they would still use their own keys to sign the action. This means it is always possible to identify which friends used the account and in what way.


This same social feature of permission mapping can be used as the infrastructure for LX copy trades. People based portfolios through these preset permission levels linking multiple account wallets to create liquidity pools that can be managed from one or several private key signatures.


\subsubsection{Permission Evaluation Applied to Copy Trading}

When delivering an action of type ``Action'' from @alice to @bob, the Lux  platform first checks to see if @alice has defined a permission mapping for @bob.groupa.subgroup.Action. If nothing is found, then a mapping for @bob.groupa.subgroup, then @bob.groupa and lastly @bob will be checked. If no further match is found, then the assumed mapping will be to the named permission group @alice.active.


Once a mapping is identified, then the signing authority is validated using the threshold multi-signature process and the authority associated with the named permission. If that fails, then it traverses up to the parent permission and ultimately to the owner permission, @alice.owner.


\subsubsection{Parallel Permission Evaluation}

The permission evaluation process is read-only and changes to permissions made by transactions do not take effect until the end of a block. This means that all keys and permission evaluation for all transactions can be executed in parallel. Furthermore, this means that a rapid validation of permission is possible without starting costly application logic that would have to be rolled back. Lastly, it means that transaction permissions can be evaluated as pending transactions and do not need to be re-evaluated as they are applied.


All things considered, permission verification represents a significant percentage of the computation required to validate transactions. Making this a read-only and trivially parallelizable process enables a dramatic increase in performance.


When replaying the blockchain to regenerate the deterministic state from the log of actions, there is no need to evaluate the permissions again. The fact that a transaction is included in a known good block is sufficient to skip this step. This dramatically reduces the computational load associated with replaying an ever growing blockchain.


\section{LX Key Capabilities}

Although BitShares was highly innovative in its conception as a product as well as its underlying technology, it has not been a widely adopted platform for crypto-asset trading. LX aims to correct the pain points of bitshares by bringing innovations into the by protocol layer to facilitate greater liquidity, upgrading the application layer to improve user experience, and injecting commercial \& operational expertise to greatly augment the adoption of a decentralized exchange. LX and the Lux Protocol are at a pinnacle moment in history to provide the financial infrastructure in this new age of global financial freedom and empowerment. Towards that end, LX and the Lux Protocol bring the following set of improvements and capabilities to bear:


\subsection{Token Curated Registry Customer Service}

The theory of customer service is based on identifying and satisfying your customers' needs and exceeding their expectations. A company must be totally committed to delivering consistently high standards of service to gain and retain customer loyalty. Everyone from top management on down must be tuned into what the customer wants. Creating a customer service culture within a company can help build success. Customer satisfaction and loyalty are inextricably linked to the quality of customer service and, ultimately, to the company's profitability.


Incentives for Content Creation


Social Anonymous Trading


Copy Trading, PBPs (people based portfolios)


\subsection{1. Atomic Swaps}

An atomic swap allows for exchange of one cryptocurrency for another without the need for a trusted third party.


For example, Alice may send BTC to Bob, while Bob sends LUX to Alice. However, as both blockchains are heterogeneous and transactions can't be undone after block consolidation, this operation introduces the counterparty risk problem - i.e.: Alice or Bob fails to honor their end of the deal.


Atomic swaps between LUX and BTC involve each party paying into a locked account. In the case of LUX, it's a multi-signature account. In the case of BTC, it's an UTXO with a scripted lock.


The creator of Litecoin, Charlie Lee, successfully implemented and demonstrated atomic swaps using LTC in exchange for BTC, Vertcoin and Decred. However, this type of swap only works between the BTC-like chains with similar scripting systems, as well as support for Check Lock Time Verification (CLTV) functionality.


With CLTV support, lock time windows can be set for the refund operation. The initiator sets a 48 hour lock time for participator payment and withdrawal. The participator sets a 24 hour lock time for initiator withdrawal. This time-locked refund scheme guarantees the atomic integrity that any party can withdraw their funds completely when the other party quits the swap process.


Because LX and the Lux Protocol don't have the BTC-like scripting system, they use a multi-signature account to lock LUX the initiator pays to the participator, and CLTV scripting to lock the BTC the participator pays to the initiator.


One drawback of this mechanism is the fact that the multi-signature method cannot implement a recovery and refund process, and LX resolves this with a requested deposit to incentivize and guarantee the trade's integrity.


Atomic Swaps				Atomic Liquidity


\subsubsection{Facilitating Liquidity}

Atomic swaps only solve the problem of trading across chains in a trustless way. It does not, however, solve the problem of how both sides of the trade find one another. In fact, the way many atomic swaps have been demonstrated technically has relied upon the fact that both trading parties do not just know of each other, but are in constant communication over some instant messaging system.


To help facilitate liquidity, LX feeds outside exchange pricing for trading pairs (e.g.: BTC vs LUX). LUX, as a coin, will be traded on different centralized exchanges after crowd-sales to all witness nodes, which will then post the fed-in price onto the chain. Users then place swap orders with desired pricing also onto the chain, the witness nodes then will automatically match these orders. During the initial bootstrapping phase of the ecosystem, Lux will incentivize experienced market makers and trading-bots to provide much needed liquidity boost. Some traditional gateways developed and operated by the LX core team or other competent ecosystem partners will long-term coexist with atomic swap capability to provide IOU service with high level security.


LX has already launched a project to integrate a LX client into a multi-cryptocurrency wallet supporting customized BTC script and multi-signature, hardware security level in second phase. This project can facilitate swaps in one wallet with one click to be more user friendly.


The swap capability between LX and other mainstream cryptocurrencies will innovate many financial tools, such as stable priced crypto-assets capable of maintaining price parity with a globally adopted currency (e.g., USD). It has high utility for convenient and censorship-resistant commerce. This can be achieved by tracking the value of a conventional underlying asset such as gold by means of an over-collateralized, counterparty risk-free, smart-contract secured blockchain loan.


\subsection{2. Exchange Traded Funds}

An ETF fund is an investment fund traded on stock exchanges, much like stocks. LX supports the following ETFs:


\subsection{Industrial ETF}

Aerospace, defense, machinery, construction, fabrication and manufacturing companies. Industry's growth is driven by demand for building construction and manufactured products like agricultural equipment.


Examples:


Industrial Select Sector SPDR (XLI A)


Vanguard Industrials ETF (VIS A+)


iShares Transportation Average ETF (IYT A)


\subsection{Financial ETF}

Investment funds, banks, real estate firms and insurance companies, among others. Revenue comes from mortgages and loans that gain value as interest rates increase.


Examples:


Financial Select Sector SPDR Fund (XLF A)


Vanguard Financials ETF (VFH A+)


SPDR S\&P Bank ETF (KBE A)


\subsection{Energy ETF}

Solar energy providers, uranium, oil and gas exploration and production companies, as well as integrated power firms, refineries and other operations. Trades occur in terms of kWh, oil gallons and other energy units.


Examples:


Energy Select Sector SPDR (XLE A)


Alerian MLP ETF (AMLP A)


Vanguard Energy ETF (VDE A)


\subsection{Healthcare ETF}

Biotechnology companies, hospital management firms, medical device manufacturers and many others. Revenues come from the constant need for human health care.


Examples:


Health Care Select Sector SPDR (XLV A)


Nasdaq Biotechnology ETF (IBB B+)


Vanguard Health Care ETF (VHT A+)


\subsection{Technology ETF}

Electronics manufacturers, software developers and information technology firms. Revenues are driven by technological upgrade cycles.


Examples:


Technology Select Sector SPDR (XLK A)


Vanguard Information Tech ETF (VGT A+)


DJ Internet Index Fund (FDN B+)


\subsection{Telecom ETF}

Wireless providers, cable companies, internet service providers and satellite companies, among others. Revenue is generated from the constant need for consumers to be connected.


Examples:


Vanguard Telecom ETF (VOX A+)


iShares US Telecommunications ETF (IYZ B+)


iShares Global Telecom ETF (IXP A-)


\subsection{Materials ETF}

Mining, refining, chemical, forestry and related companies that are focused on discovering and developing raw materials. Revenues come from raw and processed material.


Examples:


Market Vectors TR Gold Miners (GDX B+)


Materials Select Sector SPDR (XLB A)


iShares U.S. Home Construction ETF (ITB A-)


\subsection{Real Estate ETF}

Residential, industrial, and retail real estate. Revenue comes from rent income and real estate capital appreciation. Sensitive to interest rate changes.


Examples:


Vanguard REIT ETF (VNQ A+)


Vanguard Global ex-U.S. Real Estate Index Fund ETF (VNQI A+)


Schwab US REIT ETF (SCHH A)


\subsection{Utilities ETF}

Electric, gas and water companies, as well as integrated providers. Revenue comes from charging consumers and businesses that provide higher-than-average dividend yields.


Examples:


Utilities Select Sector SPDR (XLU A)


Vanguard Utilities ETF (VPU A+)


Shares Global Infrastructure ETF (IGF A-)


\subsection{Consumer Discretionary ETF}

Retailers, media companies, consumer service providers, apparel companies and consumer durables. Revenues come from an improving economy when consumer spending accelerates.


Examples:


Consumer Discretionary Select Sector SPDR (XLY A)


Consumer Discretionary AlphaDEX Fund (FDX)


Vanguard Consumer Discretion ETF (VCR A+)


\subsection{Consumer Staples ETF}

Food and beverage companies, as well as companies that create products that consumers are unwilling to cut from their budgets. In general, these companies are defensive plays capable of withstanding an economic downturn.


Examples:


Consumer Staples Select Sector SPDR (XLP A)


Consumer Staples AlphaDEX Fund (FXG B+)


Vanguard Consumer Staples ETF (VDC A+)


This approach to diversified trading aims to help investors minimize long-term risk and promote opportunities for growth.


\subsection{2. Oracles  and Crypto-asset Custody for Gateways}

Oracles are data feeds that bring data from off the blockchain (off-chain) data sources and puts it on the blockchain (on-chain) for smart contracts to use. This is necessary because smart contracts running on Ethereum cannot access information stored outside the blockchain network.


Giving smart contracts the ability to execute using off-chain data inputs extends decentralized applications' value into other real world assets such as real estate and commodities.


Lux Exchange leverages Chainlink as an Oracle to provide Verifiable Random Function (VRF)which in turn supports numerous RWA Tokens (Real World Asset Tokens) with a cryptographically secure and provably fair source of on-chain randomness (RNG), giving users a strong guarantee of application impartiality.


\subsubsection{Gateways}

Gateways are the recommended method of moving funds into or out of the LX platform. They simplify the process of moving funds from one blockchain-based crypto currency to another. Gateways are basically equivalent to the standard exchange model where the user depends on the solvency of the exchange to be able to redeem their coins.


Generally, gateways issue assets prefixed with their symbol, like LUX, APE, APD. These assets are backed 100\% by the real BTC, LUX, ETH or any other coin that people deposit with the gateways. A LUX.BTC is thus, in theory, equivalent to the BTC a user gets on Poloniex, which could be prefixed POLO.BTC. In both cases, the user relies on the service providers to remain solvent in order to back the value of the assets they've issued.


Although gateways only provide a single service, which is in itself just a single part of an exchange's overall operations, the security requirements are nevertheless high. LX, with support from its ecosystem partner, Lux Partners Ltd., has an independent crypto-asset custodian service prototype for the oracle gateways securing their holdings of crypto-assets. This critical security feature will help to streamline gateway setup and lower the risk of centralized management of high-valued crypto-assets.


\subsubsection{Hot Wallet Security}

The system maintains a percentage of digital assets in an online set of hot wallets stored within a database on the Lux  blockchain. These wallets are protected via a non-permissioned and trustless set of smart contracts designed to carry out the capability of copy trades and mirror trades as expert traders make market moves.


Any user can voluntarily withdraw their assets from the gateway, which will be automatically sent from the hot wallet, and signed with a cryptographic signature to authenticate the transaction. All seed keys to our hot wallets are produced using OTP to generate entropy and are encrypted with post quantum security hashing algorithms to ensure the safe storage of our user's digital assets.


\subsubsection{Cold Wallet}

For the gateway's remaining crypto-asset holdings, LX recommends the high-security cold wallet solution FREZOR. (coming soon! Q4 2023)


Multi-signature scheme is considered a cryptographically-secure solution for high-valued assets. With stateless script, the Bitcoin blockchain supports multi-signature from the early stage. ETH's community developed a smart contract to implement the scheme, but the security of smart contracts has been always a concern for assets holders, especially after severe vulnerabilities (e.g., Parity) were revealed during 2017. Other coins, such as Monero, have yet to release their multi-signature features. Lux Partners Ltd. provides a Multi-signature solution for multi-signers with a secret sharing algorithm leveraging post quantum cryptography. To solve the security problem that the entire private key is recovered in memory, the design uses banking level wys/wys (what you see/what you sign) hardware to share and recover the secret only inside the hardware key.


The co-signing process can be done in a decentralized manner, whether in-person or online. The entire signing process' security is not contingent upon how secure the USB hub or internet routing is by incorporating ad-hoc key algorithm to secure the P2P communications.


\subsection{3. RWA Tokens}

A RWA Token is a cryptocurrency whose value is backed by a real world asset, such as Uranium or gold. RWA tokens always have 100\% or more of their value backed by the value of the underlying asset to which they can be converted using a physical asset redemption contract at an exchange rate set by a trustworthy price feed. In all but the most extreme market conditions, RWA Tokens are guaranteed to be worth at least their face value (and perhaps more, in some circumstances).


Like any other fully defi NFT, RWA Tokens  are non- fungible, non-divisible, and free from any restrictions of transferring ownership. It is simply the certificate of ownership to the physical asset and therefore not a security. By removing the restrictions adherent to that of the transfer of ownership on traditional in-ground assets and commodities, RWA Tokens provide a solution that is hyper liquid, fully transparent and completely private leveraging the benefits of immutable ledger technology coupled with homomorphic encryption to eliminate front running on sales within the exchange.


\subsubsection{Price Stable Currencies backed by RWA Tokens}

One of the major impediments to holding any asset that is not backed by real world value is its unpredictable volatility. The more volatile an asset's relative price is to other commonly-held assets (e.g.: USD), the more risk there is in holding it. If volatility is sufficiently high, such assets become too risky to be used in everyday commercial activities. Crypto-currencies, being a relatively new asset class, is usually significantly more volatile than their non-digital counterparts. BTC/USD pricing's volatility is lowering slowly over time, but it's still around 5 to 6\%. Compare that to a more conventional asset volatility, such as Gold/USD pricing or USD/EUR exchange rate, which is only around 0.5\%.


One way to mitigate volatility and risk is to have price-stable assets in the exchange that's pegged to a relatively less-volatile asset, say, Uranium. This asset acts as a stable anchor against which all other crypto-assets could be valued, and each unit of such an asset will always give a predictable return. LX has built in price-stable assets. One such asset is a uranium-pegged stable currency that are designed to be equivalent of the TLV (total locked value) of the RWA Tokens staked into the DAO governing the rewards and incentives. These reward and incentive models mint the Price Stable currencies. The idea of such currencies is not new. BitShares for example, has such currencies for trading in its ecosystem. However, pegged currencies within the BitShares ecosystem suffer from two distinct challenges:


Pegged Currencies are in themselves too volatile (they are not backed by real world value) as it is collateralized only by BTS, a relatively volatile crypto-asset.


There are too few issuances of the pegged currencies as there are no incentives to issue them. In LX's design, the collateral is based on real world assets such as uranium and gold, a far more stable foundation of value. LX's built-in interest rates also give dynamically adjusted incentives for RWA tokens holders to issue price stable currency GLUX or ULUX.


\subsubsection{Gold as Collateral}

The functionality of a price-stable asset is that when users trade it into another asset, the trading ratio is fixed. For example, ideally, it's desirable to make sure that every single GLUX can be returned for the unstaking of the RWA Token that is staked into relative asset DAO specific to that RWA. However, given the legal and regulatory ambiguity surrounding the use of fiat currencies in a cryptocurrency exchange such as LX, LGOLD is instead traded into another asset that's more stable and liquid, such as GOLD. Hence, every LGOLD needs to be able to be traded into exactly 1 USD worth of GOLD. At the time of this writing, 1 oz of GOLD is trading at \$1,756 USD, which means 1 LGOLD needs to be traded for .00077 GOLD oz.


For someone to be able to consistently trade LGOLD into gold, there needs to be a set amount of gold ready to be traded. That is, whoever is issuing (selling) LGOLD must have a corresponding quantity of GOLD oz as collateral. To guard against the GOLD/USD kind of volatility, LX requires that the issuer of LGOLD set aside 1.25x the reserve of gold necessary to cover every single LGOLD issued. For example, to issue 1 single LGOLD, which is redeemed for 0.00077 unit of GOLD oz (equivalent to 1 USD), 0.001514 GOLD oz (equivalent to \$1.25 USD) needs to be set aside to cover any volatility.


\subsubsection{Incentives}

Having LGOLD sounds like a great incentive for someone who's holding it, but why would anyone want to issue LGOLD? That's where interest comes in. For every unit of LGOLD issued, the holder of the asset needs to pay the issuer of the asset an interest rate. Think of it as a fee paid to someone who's providing this big pool of collateral as a cushion against market volatilities. Since LX is a free-market type of decentralized exchange, it isn't guaranteed that 1 LGOLD will always translate into exactly 1 USD worth of GOLD. When this happens, LX adjusts interest rates to make sure that issuers are properly incentivized to issue more or less of LGOLD in order to balance out supply and demand. For example, if demand for LGOLD is rising, that means 1 LGOLD will trade for more than 1 USD worth of GOLD. When this happens, it is desirable to increase the supply of LGOLD to bring that ratio back to a 1 to 1 trading ratio by increasing the interest rate. When the reverse happens, it's desirable to lower the interest rate to reduce the supply.


\subsubsection{Interest Rate}

LX uses LGOLD vs USD real time premium to generate the current (per 8 hour) interest rate. The annualized interest rate moves from 0\\% to a 15\\% cap, while the premium goes from 0\% up.


In LX initial iteration, the interest rate is calculated in the following manner:


b\_LUX = \$ GOLD oz value traded for 1 LGOLD on LX


b\_market = \$ GOLD oz value being traded for 1 USD on a collection of major exchanges


premium = (b\_LUX - b\_market)/b\_market


Setting the interest rate:


if (premium = 0\%): interest\_rate = 0\%


if (0\% < premium ≤ 2\%): interest\_rate = 2\%


if (2\% < premium ≤ 15\%): interest\_rate = premium


if (premium > 15\%): interest\_rate = 15\%


The initial 2\% is set to create a threshold to help trigger action on the part of the issuer - people tend to understand a fixed threshold better and are more likely to act upon one. LX also caps the interest rate at 15\% to make sure that the holder of the asset won't feel insecure about a potentially limitless interest rate that continuously rises.


The system then automatically calculates the interest rate based on the average premium over the past 8 hours on LX and then takes a snapshot of all LGOLD holding and issuances at the end of that 8 hour period. Once every 8 hours, the holders of LGOLD's account is deducted by the interest (they will have less LGOLD) and the issuers' accounts receives deposits of LGOLD by the exact same amount. All deductions and deposits are done in proportion to the snapshot of LGOLD holdings and issuances.


\section{Development Roadmap}

The LX development roadmap is divided in 7 phases:

\begin{enumerate}
  \item LUX Constitution
  \item Lux DAO
  \item Lux Protocol
  \item DEX Core Platform
  \item UI/UX
  \item DApp
  \item Hardware Wallet Integration
\end{enumerate}

\subsection{1. LUX Constitution and Ricardian Contracts}

The LUX Constitution and Ricardian Contracts are written with help of specialized legal and technical consultancy. Any relative and corresponding regulatory compliance and licensing required to interact within particular jurisdictions will be filed and maintained for optimum efficiency to operate with both retail and institutional investors globally.


\subsection{2. Lux DAO}

The Lux DAO leverages holographic consensus, quadratic voting, aged staking and weighted POS based on experience points gained through reputation and verified with proof of human biometric digital signal processing to authenticate user identity.


\subsection{3. Lux Protocol}

The Lux Protocol implementation depends on the following items:


LX Constitution


LX Token


LX Permission Mapping


LX Smart Contracts


C++, CMake, ABI, and Ricardian Contracts for the following LUX Contracts:


Exchange Contract


Multi Index Containers Contract


Identity Contract


Order Book Contract


Decentralized Liquidity Pool Contract


Atomic Swap Contract


Gateway Contract


Permission Mapping Contract


Token Curated Registry Contract


ETF Contract


Copy Trading Contract


PBP Contract


LX Plugins


C++, CMake and ABI for the following LX Plugins:


Lux Network Plugin


Lux Producer Plugin


Lux Order Book Plugin


Lux Decentralized Liquidity Pool Plugin


Lux Gateway Plugin


LX API


C++ API with the following functions:


Deposit


Check Balance


Withdraw


Create Order


Cancel Order


Execute Order


View Order Book


\subsection{3. DEX Core Platform}

The LX DEX core platform will be carefully designed and implemented with Quality Assurance of all C++ code will enforce IEC 61508 and CERT compliance.


The DEX core functionality will be written in Solidity. The platform consists of the following list of Lux  Smart Contracts and Node Plugins:


LX DEX:


Order Book Contract


Atomic Swap Contract


Decentralized Liquidity Pool Contract


Permissioned Mapping Contract


Gateway Contract


Hot Wallets


Cold Wallets


Exchange Traded Fund Contract


Market Pegged Asset Contract


LUX Blockchain Infrastructure


LUX Chain Node


LUX Producer Node


LUX Order Book Node


LUX Decentralized Liquidity Pool Node


LUX Gateway Node


Cloudflare


\subsection{4 DApp UI/UX}

The LX DApp User Interface (UI) and User Experience (UX) must be carefully designed and implemented with the use of lean startup customer development. The LX User Interface must be written in JavaScript with desktop and mobile support for Google Chrome, Mozilla Firefox, Apple Safari, and Microsoft Edge.


The DApp UI/UX must provide of the following items:


Account Management


SHA3 authentication with email and high entropy password


Yubikey NEO for U2F


Wallet Management


Integration with common wallet providers


Lux Wallet


Key Management


DApp keychain protected by high entropy password


Yubikey NEO forgot optional safe key storage


People Based Portfolios (Lux  Social Contracts)


Copy Trading (Lux  Social Contracts)


Messaging (Engine Framework)


Token Curated Registries


Customer Support


Use the C++ Engine Framework to write dApp chat sessions with token rewards proportional to user feedback about the representative's performance


Rewarded Markdown Blogging


GitHub integration


LUX Token Curated Rewarding proportional to upvote and downvote of each post


Write posts for development of LUX Trading Bots in C++


\subsection{5 Hardware Wallet Integration}

The LX DApp must support the Yubikey NEO hardware wallet.


\subsection{6 Quality Assurance}

All code must be tested by QA-SYSTEM's QA-C and QA-C++ tools to perform CERT and IEC 61508 Quality Assurance.


The IEC61508 (Functional Safety of Electrical/Electronic/Programmable Electronic Safety-related Systems), is an international standard widely used in Industrial Automation. The IEC61508 standard also requires use of coding standards, such as MISRA and CERT.


CERT is maintained by the Software Engineering Institute (SEI), a research and development center primarily funded by the U.S. Department of Defense and the Department of Homeland Security. The CERT Division at SEI is operated by Carnegie Mellon University and responsible for publishing these standards. The CERT\textregistered{} Secure Coding Standards for C and C++ are standards that provide rules and recommendations that target insecure coding practices and undefined behaviors that can lead to exploitable vulnerabilities.


It is our goal to maintain IEC 61508 and CERT compliance in all code design and development. There are many existing and emerging standards that might be applied but the IEC 61508 and CERT specifications align closely with our commitment to provide quality services and error free processing. It achieve this goal the C++ code must be tested to assure that:


Quality Assurance of all C++ code must be ensured for IEC 61508 and CERT compliance.he


The C++ code must also be tested against:


Application level vulnerabilities


Cross domain information leakage (Onchain name servers for domain registry)


Client side logic and data storage vulnerabilities


Cloud configurations vulnerabilities


Blockchain level vulnerabilities


Input validation and representation


API abuse


Access control mechanisms


Memory management


Time and state


Error and exception handling


Encapsulation and hidden defects


Flaws at browser runtime


Insecure JNI


Unused variables after assignment


\end{document}
